% page 111 of the facsimile (image p-110 of the scan)
% block 175.3 x 229.0 mm ; left margin 26.7 mm
\pgc{5.080mm}{0.000mm}{
\l{\cel{55}lO3}
\l{}
\l{}
\l{\sou{dento}. I. (\sou{Anat.)}\cel{3}Singla ek la mikra osti emalioza qui garnisas}
\l{\cel{3}la maxilo e la mandibulo di la homo, di ula animali, ed uzesas}
\l{\cel{3}da li por tranchar, arachar, pistar la nutrivi, por mordar, e c.}
\l{\cel{3}- II. (\sou{Tekniko})\cel{2}To quo esas dentatra. - EFIRS.}
\l{}
\l{\sou{dentelo}. Texuro ajuroza, ek reto tenua de filo o de silko, di}
\l{\cel{3}qua la bordo esas generale dentoza, uzata kom fundo por flori,}
\l{\cel{3}ornivi facita per spindelo o per agulo. - F.}
\l{}
\l{\sou{dentifrico}. Preparajo qua uzesas por frotar, por netigar la denti.}
\l{\cel{3}- EFIS.}
\l{}
\l{\sou{denuncar}. (\sou{trans.}) Indikar (ulu) kom kulpinta, precipue a la tribu-}
\l{\cel{3}nali od a la policisti. - DEFIRS.}
\l{}
\l{\sou{departar}. (\sou{netrans., de,}\cel{1}ad). Forirar de loko. - EFIS.}
\l{}
\l{\sou{departmento}. Dividuro teritoriala, o resortiso teritoriala. - DEFIRS.}
\l{}
\l{\sou{dependar}. (\sou{netrans.,}\cel{1}de). Ligesar necese a persono od a kozo, kom}
\l{\cel{3}nepovanta esor o realigesor sen lu. - EFIS.}
\l{}
\l{\sou{deperisar}. (\sou{netrans}.)Tale febleskar ke morto probable eventos balde.}
\l{\cel{3}- FI.}
\l{}
\l{\sou{depesho}. Informo (publika o privata) sendita (ad ulu) per moyeno}
\l{\cel{3}rapida. - DEFIRS.}
\l{}
\l{\sou{deplorar}.\cel{2}(\sou{trans.}) Manifestar, expresar, profunda chagreno pro}
\l{\cel{2}(ulo). EFIRS.}
\l{}
\l{\sou{deponenta}. (\sou{Gram. Latina}). Qua havas la formo pasiva e la senco}
\l{\cel{3}aktiva. - DEFIS.}
\l{}
\l{\sou{deportar}. (\sou{trans}.)Duktar (kondamnito) adexter la teritorio ad-an}
\l{\cel{3}loko de ube lu devos ne ekirar. - DEFIRS.}
\l{}
\l{\sou{depozar}. (\sou{trans}.) Pozar (kozo) en loko ube lu esos sekura. - DEFIRS.}
\l{}
\l{\sou{depresar}. (\sou{trans.}) Abasar la nivelo per preso de alte ad infre.}
\l{\cel{3}- DEFIS.}
\l{}
\l{\sou{deputar}. (\sou{trans.})\cel{2}Elektar (ulu) por reprezentar la naciono en}
\l{\cel{3}asemblitaro deliberanta. - DEFIRS.}
\l{}
\l{\sou{derivar}. (\sou{trans}.)\cel{2}I.Ekirigar de lua lito aquo fluanta, por direktar}
\l{\cel{3}lu ad altra direciono. II. Atraktar (sango, humoro) de parto di}
\l{\cel{3}la korpo ad altra, ube lia adfluo esas min danjeroza. III. Ek-}
\l{\cel{3}irigar (vorto) de altra per modifikar lua formo ; (senco) de}
\l{\cel{3}altra senco per extenso, e c. - EFIS.}
}
